\section{Applications}
\label{sec:applications}

The measurements of \Cref{sec:eval} explain where the AVL tree earns its place
and where it does not. It keeps the shortest tree of any of the structures in
\Cref{tab:comparison}, and it pays for that during updates. A workload therefore
suits an AVL tree when it searches far more often than it modifies, and when the
data is in memory so that a comparison, rather than a block fetch, is the unit of
cost.

\subsection{In-memory ordered collections}

The clearest case is an ordered dictionary that is read constantly and written
occasionally. A routing table, a symbol table in a compiler, or an index over a
data set loaded once and queried many times all have this shape. Each search saves
comparisons against a red-black tree holding the same keys, and the extra
rotation work is charged only to the rare update.

The advantage over a hash table is the ordering. A hash table answers an exact
lookup in expected constant time and beats any search tree at that one operation,
but it cannot answer a range query, report the keys in sorted order, or find the
key nearest a target without scanning everything. Those operations are what a
search tree provides, and they are the reason a program picks one.

\subsection{Guaranteed worst-case behaviour}

The second case is a system that needs a bound on every individual operation
rather than on a sequence of them. \Cref{thm:heightbound} bounds the height of
the tree at all times, so \Cref{cor:opcost} bounds every search, and
\Cref{cor:insertcost} and \Cref{thm:deletebound} bound every update. No
particular operation can be slow.

This distinguishes the AVL tree from both of its natural competitors for
in-memory work. A hash table has an expected-case guarantee that a bad set of
keys can defeat, and a rehash pauses the structure for a time proportional to its
size. A splay tree has an amortised guarantee, so a single access may cost
$\Theta(n)$ even though a long run of accesses is cheap. Where a deadline applies
to each operation, as in scheduling or control software, neither of those bounds
is usable and the worst-case bound is.

\subsection{Where other structures are chosen instead}

It is worth being clear about the limits, because the AVL tree is often not the
right choice.

Workloads dominated by insertions and deletions favour the red-black tree, whose
deletion needs at most three rotations against the $O(\log n)$ of
\Cref{thm:deletebound}. This is why the ordered maps in several widely used
standard libraries, and the kernel data structures that track process and memory
ranges, are red-black rather than AVL.

Data that lives on disk or on a solid-state drive favours a B-tree, for the
reason given at the end of \Cref{sec:comparison}: the cost is the number of
blocks fetched, and a shallow, wide tree fetches fewer. Database indexes are
B-trees for this reason and not because their designers were unaware of the
height bound.

Workloads with strong locality of reference favour a splay tree, which moves each
accessed key to the root and so serves a small working set from the top of the
tree. The AVL tree treats every key alike and cannot exploit that pattern.

Finally, a workload that only ever needs exact lookup, with no ordering
requirement, should use a hash table. The AVL tree's ordering guarantees cost
something to maintain, and a program that does not use them should not pay for
them.
